\documentclass[11pt, a4paper, twoside]{article}

% --- Packages ---
\usepackage[utf8]{inputenc}
\usepackage{amsmath, amssymb, amsthm}
\usepackage{geometry}
\usepackage{tcolorbox}
\usepackage{tikz-cd}
\usepackage{fancyhdr}

% --- Page Layout & Styling ---
\geometry{
	top=1in, 
	bottom=1in, 
	left=1.25in, 
	right=1.25in
}

\pagestyle{fancy}
\fancyhf{}
\fancyhead[LE,RO]{\thepage}
\fancyhead[LO]{\small \textit{Lecture Notes: Linear Block Codes}}
\fancyhead[RE]{\small \textit{The Coding Theory of Cohomology}}
\renewcommand{\headrulewidth}{0.4pt}

% Custom boxes
\newtcolorbox{conceptbox}[1][]{
	colback=purple!5!white, colframe=purple!75!black,
	fonttitle=\bfseries, title=#1, arc=3pt, boxrule=1pt,
	left=10pt, right=10pt, top=10pt, bottom=10pt
}

% --- Macros ---
\newcommand{\F}{\mathbb{F}}
\newcommand{\im}{\text{im}}
\newcommand{\kerop}{\text{ker}}
\newcommand{\cH}{\mathcal{H}}

\begin{document}
	
	\begin{center}
		{\Large \textbf{The Coding Theoretical Meaning of Cohomology}} \\[0.5em]
		{\normalsize \textit{From Classical Exactness to Topological Quantum Error Correction}}
	\end{center}
	\vspace{1em}
	
	We define a coding chain complex where $H \circ G = 0$, but we deliberately relax the exactness condition at the middle node. This allows $\im(G) \subsetneq \kerop(H)$. 
	
	The failure of this sequence to be exact is measured by the first cohomology group:
	$$ \cH^1 = \frac{\kerop(H)}{\im(G)} $$
	
	In classical coding, a non-zero cohomology group means the code is broken (there exist valid syndromes that do not correspond to any generated message). However, in \textbf{Topological Quantum Error Correction} (such as CSS codes), the cohomology group is the precise location where protected information is stored.
	
	\vspace{1em}
	
	\begin{conceptbox}[1. The Trivial Fluctuations: $\im(G)$]
		In a quantum topological code, the image of $G$ no longer represents encoded messages. Instead, $\im(G)$ represents the \textbf{Stabilizer Group} (or gauge transformations). 
		
		These are physical operations or local errors that alter the state of the physical hardware but \textit{do not change the logical information}. If an error $e \in \im(G)$ occurs, the system's physical configuration changes, but the data remains perfectly intact. Geometrically, these correspond to ``exact forms'' or contractible boundaries.
	\end{conceptbox}
	
	\vspace{1em}
	
	\begin{conceptbox}[2. The Code Space: $\kerop(H)$]
		The kernel of the parity-check matrix $H$ consists of all physical states that yield a syndrome of zero ($s = \mathbf{0}$). To the error-correcting hardware, any state in $\kerop(H)$ appears to be a valid, error-free configuration.
	\end{conceptbox}
	
	\vspace{1em}
	
%	\begin{conceptbox}[3. The Logical Qubits: $\cH^1 = \kerop(H) / \im(G)$]
		Because the space is not exact, there are states that look valid to the hardware but were not generated by local stabilizers. The \textbf{Logical Information} is stored in the equivalence classes of the quotient space $\cH^1$.
		
		By modding out $\im(G)$, the mathematics states: \textit{``Two valid physical configurations encode the exact same logical information if and only if they differ by a harmless stabilizer operation.''} 
		The number of protected logical qubits $k_{log}$ is exactly the dimension of this cohomology group:
		$$ k_{log} = \dim \cH^1 $$
%	\end{conceptbox}
	
	\vspace{1em}
	
	\begin{conceptbox}[4. Cohomological Errors: Logical Failure]
		What happens if the environment induces an error $e \in \F_q^n$ such that $e$ is closed ($e \in \kerop(H)$) but not exact ($e \notin \im(G)$)?
		\begin{itemize}
			\item Because $e \in \kerop(H)$, the hardware computes $s = He^T = \mathbf{0}$. The hardware believes no error has occurred.
			\item Because $e \notin \im(G)$, the error is not a harmless gauge transformation.
		\end{itemize}
		This error effectively shifts the system from one cohomology class to another. This represents an \textbf{undetectable logical error}. Geometrically, this corresponds to an error chain that wraps entirely around a topological ``hole'' in the manifold (e.g., a non-contractible loop on a torus). The design of a topological code aims to make the macroscopic distance required to complete such a cohomological error physically improbable.
	\end{conceptbox}
	
\end{document}